%=====================================================================
% Priorización de requisitos y casos de uso
%   Las tablas de esta sección y las columnas Categoría y Prioridad de
%   lista-casos.tex salen de los mismos datos; si cambia una, cambiar
%   la otra, y también la fila "Prioridad" de cada caso en casos/.
%=====================================================================
\section*{Priorización de requisitos y casos de uso}
\addcontentsline{toc}{section}{Priorización de requisitos y casos de uso}

El proyecto tiene catorce semanas y debe entregar cada semana un incremento que pueda auditarse. Los 49 casos de uso no caben en ese plazo, así que esta sección decide cuáles se construyen primero. Ningún caso se elimina: los que no entran se marcan como \textbf{fase posterior} y conservan su redacción completa, para retomarlos sin rehacer el análisis.

\subsection*{Requisitos obligatorios del cliente}
\addcontentsline{toc}{subsection}{Requisitos obligatorios del cliente}

Estas restricciones las fija el cliente y no son negociables: el sistema no se entrega sin cumplirlas todas, y el resto del alcance se prioriza alrededor de ellas. Los casos \mbox{CU-01} y \mbox{CU-02} las citan en su trazabilidad con identificadores \mbox{CON-xx}; las que no tienen identificador las confirmó el equipo.

\begin{center}\small
\begin{longtable}{|L{0.25\textwidth}|L{0.14\textwidth}|L{0.49\textwidth}|}
\hline
\textbf{Restricción obligatoria} & \textbf{Identificador} & \textbf{Dónde se cumple} \\ \hline
\endfirsthead
\hline
\textbf{Restricción obligatoria} & \textbf{Identificador} & \textbf{Dónde se cumple} \\ \hline
\endhead
Programar en C\# & \mbox{CON-01} & Todo el sistema: cliente y Servidor de partidas. \\ \hline
Conexión por TCP, sin WebSockets & \mbox{CON-02}, \mbox{CON-03} & Todo mensaje entre el cliente y el Servidor de partidas, como \texttt{MOVE\_\allowbreak{}PAWN\_\allowbreak{}REQUEST} del \mbox{CU-20}. \\ \hline
Base de datos en SQL Server & \mbox{CON-04} & Modelo de datos e implementación. \\ \hline
Multijugador en línea & sin identificador en el repositorio & \mbox{CU-17}, \mbox{CU-18}, \mbox{CU-19}, \mbox{CU-20}, \mbox{CU-21}, \mbox{CU-24}, \mbox{CU-25} y \mbox{CU-26}. \\ \hline
Segundo factor de autenticación & \mbox{CON-08} & \mbox{CU-01} lo pide al iniciar sesión; \mbox{CU-09} \mbox{FA-10} lo activa (\mbox{Q-03}). \\ \hline
Moderación de chat & \mbox{CON-09} & \mbox{CU-28} aplica el filtro de palabras; \mbox{CU-33}, \mbox{CU-42} a \mbox{CU-45} reportan, revisan, sancionan y apelan. \\ \hline
Catorce semanas con incrementos semanales auditables & sin identificador en el repositorio & Esta priorización: los casos \textit{Must} caben en el plazo y cada incremento cierra casos completos. \\ \hline
\end{longtable}
\end{center}

Los casos citan además otras restricciones cuyo origen no pudo confirmarse: \mbox{CON-05} (texto Unicode), \mbox{CON-07} (respuesta en menos de un segundo), \mbox{CON-11} (cambio de idioma) y \mbox{CON-12} (comprensible para un niño de ocho años). Se tratan como restricciones del proyecto y se siguen cumpliendo en los casos que entran. Ningún caso cita \mbox{CON-06} ni \mbox{CON-10}; conviene recuperar su texto del documento de requisitos del cliente.

\subsection*{Criterio}
\addcontentsline{toc}{subsection}{Criterio de priorización}

Cada caso recibe una \textbf{categoría}, que dice de dónde sale, y una \textbf{prioridad} con el método MoSCoW.

\begin{reglas}
  \item \textbf{(a) Obligatorio.} Requisito del cliente, no negociable. El caso existe porque lo exige una restricción de la tabla anterior, o porque otro caso obligatorio no funciona sin él: el \mbox{CU-04} entra porque el \mbox{CU-01} rechaza las cuentas sin verificar.
  \item \textbf{(b) Núcleo de jugabilidad.} Las reglas de Quoridor y lo mínimo para jugarlas en línea: mover el peón, el salto, colocar muros sin encerrar a nadie, turnos y reloj, condición de victoria, partidas de 2 y 4 jugadores, buscar partida, sala privada, reconexión y fin de partida.
  \item \textbf{(c) Agregado del equipo.} Todo lo demás: skins, cosméticos, tienda, cajas, monedas, niveles, divisiones, ranking, historial, repeticiones, espectador, emotes, perfil y enlaces de perfil, amigos, invitado, IA y tutorial.
\end{reglas}

\begin{reglas}
  \item \textbf{Must.} Debe estar en la entrega. Sin él no se cumple una restricción obligatoria del cliente o no se puede jugar una partida completa.
  \item \textbf{Should.} Debería estar. Se construye en cuanto los \textit{Must} estén cerrados.
  \item \textbf{Could.} Podría estar. Es lo primero que se recorta si el plazo aprieta.
  \item \textbf{Won't (fase posterior).} No entra en las catorce semanas. Por regla, todo agregado del equipo cae aquí.
\end{reglas}

\textbf{La jugabilidad va primero.} Dentro de los \textit{Must}, los casos de jugabilidad preceden a los obligatorios del cliente: una partida completa entre dos clientes por TCP es el primer incremento que demuestra algo, y el inicio de sesión con segundo factor, el chat moderado y los reportes se montan encima. Una partida nunca tiene tres jugadores: los modos admiten solo 2 o 4 (\mbox{CU-17} \mbox{RN-02}), y el modelo de datos lo fija con un \texttt{CHECK}.

\subsection*{Núcleo de jugabilidad}
\addcontentsline{toc}{subsection}{Núcleo de jugabilidad}

Cada regla del núcleo ya está especificada en un caso de uso. La tabla dice dónde, para que el primer incremento se construya y se pruebe contra esas reglas.

\begin{center}\small
\begin{longtable}{|L{0.28\textwidth}|L{0.21\textwidth}|L{0.39\textwidth}|}
\hline
\textbf{Regla o función} & \textbf{Dónde está} & \textbf{Qué fija} \\ \hline
\endfirsthead
\hline
\textbf{Regla o función} & \textbf{Dónde está} & \textbf{Qué fija} \\ \hline
\endhead
Mover el peón & \mbox{CU-20} \mbox{RN-02}, \mbox{RN-04} & Una casilla en ortogonal, libre y sin muro de por medio. \\ \hline
Salto recto y diagonal & \mbox{CU-20} \mbox{RN-03} & Recto sobre un rival adyacente; en diagonal si detrás hay muro u otro peón. \\ \hline
Colocar muros sin encerrar a nadie & \mbox{CU-21} \mbox{RN-02} a \mbox{RN-05} & Muros por modo, de dos casillas, sin solaparse, y nunca sin camino a la meta. \\ \hline
Turnos & \mbox{CU-20} \mbox{RN-05}, \mbox{RN-06}; \mbox{CU-21} \mbox{RN-06} & Una sola acción por turno y jugadas numeradas. \\ \hline
Reloj & \mbox{CU-20} \mbox{FA-06}, \mbox{CU-25} \mbox{RN-01} & El servidor descuenta el tiempo; agotarlo termina la partida. \\ \hline
Condición de victoria & \mbox{CU-25} \mbox{RN-01}, \mbox{RN-02} & Llegar a la fila de meta; en cuatro jugadores se clasifica y la partida sigue. \\ \hline
Partidas de 2 y 4 jugadores & \mbox{CU-17} \mbox{RN-02}, \mbox{CU-18} & Clásico 9×9 y rápida 7×7 con dos jugadores; 9×9 con cuatro. \\ \hline
Buscar partida & \mbox{CU-17} & Cola por modo y reloj con ventana de elo. \\ \hline
Sala privada & \mbox{CU-18}, \mbox{CU-19} & Sala con código de seis caracteres y ajustes del anfitrión. \\ \hline
Reconexión & \mbox{CU-26}, \mbox{CU-24} & La partida espera con el reloj corriendo; a los sesenta segundos el rival puede reclamarla (\mbox{D-15}). \\ \hline
Fin de partida & \mbox{CU-25}, \mbox{CU-23}, \mbox{CU-24} & Meta, tiempo agotado, rendición, tablas o abandono. \\ \hline
El servidor manda & \mbox{CU-20} \mbox{RN-01}, \mbox{RN-07}; \mbox{CU-21} \mbox{RN-08} & El cliente propone y el servidor valida cada jugada sobre su propio estado. \\ \hline
\end{longtable}
\end{center}

\subsection*{Prioridad de cada caso de uso}
\addcontentsline{toc}{subsection}{Prioridad de cada caso de uso}

La misma prioridad aparece en la lista de casos de uso y en la fila \textit{Prioridad} de la ficha de cada caso.

\begin{center}\small
\begin{longtable}{|L{0.075\textwidth}|L{0.25\textwidth}|L{0.11\textwidth}|L{0.44\textwidth}|}
\hline
\textbf{ID} & \textbf{Caso de uso} & \textbf{Categoría} & \textbf{Por qué} \\ \hline
\endfirsthead
\hline
\textbf{ID} & \textbf{Caso de uso} & \textbf{Categoría} & \textbf{Por qué} \\ \hline
\endhead
\multicolumn{4}{|l|}{\textbf{Must: debe estar en la entrega (17 casos)}} \\ \hline
\mbox{CU-17} & Elegir modo y buscar partida clasificatoria & Jugabilidad & Buscar partida en línea de 2 o 4 jugadores; cubre también el multijugador en línea. \\ \hline
\mbox{CU-18} & Crear una sala privada & Jugabilidad & Sala privada: partida en línea entre conocidos, sin depender de la cola. \\ \hline
\mbox{CU-19} & Unirse a una sala con código & Jugabilidad & Entrar a la sala privada con su código. \\ \hline
\mbox{CU-20} & Mover el peón & Jugabilidad & Mover el peón y el salto recto o diagonal: regla central de Quoridor, con turno y reloj. \\ \hline
\mbox{CU-21} & Colocar un muro & Jugabilidad & Colocar muros sin encerrar a nadie: la otra regla central. \\ \hline
\mbox{CU-24} & Abandonar la partida & Jugabilidad & Abandono y reclamo de victoria: sin él una partida con un jugador caído no termina. \\ \hline
\mbox{CU-25} & Finalizar la partida y aplicar la progresión & Jugabilidad & Condición de victoria, tiempo agotado y cierre de la partida con su elo. \\ \hline
\mbox{CU-26} & Reconectar a una partida en curso & Jugabilidad & Reconexión: una desconexión no debe costar la partida. \\ \hline
\mbox{CU-01} & Iniciar sesión & Obligatorio & Inicio de sesión con segundo factor (\mbox{CON-08}); sin sesión no hay juego en línea. \\ \hline
\mbox{CU-02} & Registrar cuenta & Obligatorio & Sin cuenta registrada no hay inicio de sesión ni segundo factor. \\ \hline
\mbox{CU-04} & Verificar correo y activar la cuenta & Obligatorio & El \mbox{CU-01} rechaza las cuentas \texttt{PENDIENTE}: sin verificar el correo nadie entra (Matriz 2). \\ \hline
\mbox{CU-05} & Cerrar sesión & Obligatorio & Cierra la sesión que abre el \mbox{CU-01}; es el mínimo del acceso. \\ \hline
\mbox{CU-09} & Cambiar contraseña o correo & Obligatorio & Su \mbox{FA-10} activa el segundo factor (\mbox{Q-03}); cambiar contraseña y correo lo acompaña. \\ \hline
\mbox{CU-28} & Enviar un mensaje de chat & Obligatorio & Moderación de chat (\mbox{CON-09}): el filtro de palabras se aplica aquí. \\ \hline
\mbox{CU-42} & Reportar a un jugador & Obligatorio & Reportar (reportante y reportado): entrada de la moderación. \\ \hline
\mbox{CU-43} & Revisar un reporte y resolverlo & Obligatorio & Revisar y resolver reportes: la moderación propiamente dicha. \\ \hline
\mbox{CU-44} & Aplicar una sanción & Obligatorio & Aplicar la sanción de cuenta o de chat. \\ \hline
\multicolumn{4}{|l|}{\textbf{Should: debería estar (5 casos)}} \\ \hline
\mbox{CU-23} & Rendirse & Jugabilidad & Rendirse da una salida limpia sin forzar un abandono. \\ \hline
\mbox{CU-03} & Recuperar contraseña & Obligatorio & Complementa el acceso: sin él una contraseña olvidada pierde la cuenta, pero el juego funciona. \\ \hline
\mbox{CU-33} & Silenciar o bloquear a un jugador & Obligatorio & Silenciar y bloquear son la moderación que ejerce el propio jugador. \\ \hline
\mbox{CU-45} & Apelar una sanción & Obligatorio & Apelar corrige errores de moderación; sin él un moderador decide sin revisión. \\ \hline
\mbox{CU-46} & Agregar un moderador & Obligatorio & Sin él los moderadores se dan de alta por SQL, como los administradores (\mbox{D-09}). \\ \hline
\multicolumn{4}{|l|}{\textbf{Could: podría estar (4 casos)}} \\ \hline
\mbox{CU-22} & Ofrecer y aceptar tablas & Jugabilidad & Las tablas no son regla de Quoridor; útiles pero prescindibles. \\ \hline
\mbox{CU-10} & Gestionar sesiones activas & Obligatorio & Cerrar sesiones remotas es seguridad deseable, no exigida. \\ \hline
\mbox{CU-47} & Retirar a un moderador & Obligatorio & Retirar el rol puede hacerse por SQL mientras tanto. \\ \hline
\mbox{CU-48} & Consultar las bitácoras & Obligatorio & Las bitácoras se escriben siempre; consultarlas desde la aplicación es deseable. \\ \hline
\multicolumn{4}{|l|}{\textbf{Won't en esta fase: fase posterior (23 casos)}} \\ \hline
\mbox{CU-06} & Jugar como invitado & Agregado & Jugar sin cuenta esquiva el segundo factor, que es obligatorio. \\ \hline
\mbox{CU-07} & Vincular la cuenta de invitado & Agregado & Solo existe si existe el invitado (\mbox{CU-06}). \\ \hline
\mbox{CU-08} & Configurar la cuenta la primera vez & Agregado & Elegir peón e icono es personalización; el peón predeterminado basta. \\ \hline
\mbox{CU-11} & Borrar la cuenta & Agregado & La baja de cuenta no es un requisito obligatorio del cliente. \\ \hline
\mbox{CU-12} & Editar el perfil: foto, icono, título y enlaces & Agregado & Foto, icono, título y enlaces de perfil son agregados. \\ \hline
\mbox{CU-13} & Cambiar el nombre (una sola vez) & Agregado & Cambio de nombre: agregado; el nombre se elige en el registro. \\ \hline
\mbox{CU-14} & Equipar un objeto cosmético en una ranura & Agregado & Cosméticos y skins. \\ \hline
\mbox{CU-15} & Consultar el perfil propio & Agregado & El perfil con estadísticas es un agregado. \\ \hline
\mbox{CU-16} & Consultar el perfil de otro jugador & Agregado & El perfil de otro jugador es un agregado; se reporta desde la partida o el chat. \\ \hline
\mbox{CU-27} & Jugar contra la IA & Agregado & La IA es un agregado y exige un algoritmo propio. \\ \hline
\mbox{CU-29} & Enviar un emote & Agregado & Emotes. \\ \hline
\mbox{CU-30} & Agregar a un amigo & Agregado & Amigos: agregado del equipo. \\ \hline
\mbox{CU-31} & Responder una solicitud de amistad & Agregado & Amigos: agregado del equipo. \\ \hline
\mbox{CU-32} & Invitar a un amigo a jugar & Agregado & Invitar amigos depende de la lista de amigos. \\ \hline
\mbox{CU-34} & Ver una partida como espectador & Agregado & Espectador. \\ \hline
\mbox{CU-35} & Consultar el ranking & Agregado & El ranking visible es un agregado; el elo sí se calcula para emparejar. \\ \hline
\mbox{CU-36} & Consultar el historial de partidas & Agregado & Historial de partidas: las partidas ya se guardan, falta solo la pantalla. \\ \hline
\mbox{CU-37} & Ver la repetición de una partida & Agregado & Repeticiones. \\ \hline
\mbox{CU-38} & Comprar un objeto en la tienda & Agregado & Tienda y monedas. \\ \hline
\mbox{CU-39} & Comprar y abrir una caja & Agregado & Cajas. \\ \hline
\mbox{CU-40} & Consultar el historial de monedas & Agregado & Monedas. \\ \hline
\mbox{CU-41} & Avanzar en el tutorial & Agregado & Tutorial interactivo; mientras tanto basta una pantalla de reglas. \\ \hline
\mbox{CU-49} & Eliminar a un amigo & Agregado & Amigos: agregado del equipo. \\ \hline
\end{longtable}
\end{center}

\textbf{Excepción dentro de un agregado.} Del \mbox{CU-12} se adelanta solo el cambio de idioma (\mbox{CON-11}, \mbox{D-21}): el idioma se elige en el registro y se cambia en los ajustes sin esperar al resto del perfil.

\subsection*{Partes de casos del núcleo que se difieren}
\addcontentsline{toc}{subsection}{Partes de casos del núcleo que se difieren}

Algunos casos que sí entran tocan estructuras de los agregados: el alta otorga objetos cosméticos, el cierre de partida reparte monedas y cajas. Esos pasos se omiten hasta la fase posterior; el resto del caso se construye tal como está redactado. Las tablas que desaparecen del esquema del núcleo por este motivo están en la sección \textit{Modelo de datos}.

\begin{center}\small
\begin{longtable}{|L{0.11\textwidth}|L{0.12\textwidth}|L{0.64\textwidth}|}
\hline
\textbf{Caso} & \textbf{Prioridad} & \textbf{Qué se difiere a la fase posterior} \\ \hline
\endfirsthead
\hline
\textbf{Caso} & \textbf{Prioridad} & \textbf{Qué se difiere a la fase posterior} \\ \hline
\endhead
\mbox{CU-01} & Must & Leer el icono y el equipamiento para el menú. \\ \hline
\mbox{CU-02} & Must & Otorgar objetos iniciales y equiparlos (\textit{UsuarioObjeto}, \textit{Equipamiento}). \\ \hline
\mbox{CU-17} & Must & Congelar el aspecto (\textit{ParticipacionObjeto}). \\ \hline
\mbox{CU-18} & Must & Congelar el aspecto (\textit{ParticipacionObjeto}). \\ \hline
\mbox{CU-25} & Must & Experiencia, nivel, monedas, cajas, divisiones y enlaces de espectador (pasos 11 y 12, \mbox{FA-02}, \mbox{FA-03}). \\ \hline
\mbox{CU-28} & Must & Canales \texttt{GLOBAL}, \texttt{AMIGOS} y \texttt{ESPECTADORES}. \\ \hline
\mbox{CU-33} & Should & Borrar amistad, solicitudes e invitaciones del par. \\ \hline
\end{longtable}
\end{center}

\subsection*{Resumen}
\addcontentsline{toc}{subsection}{Resumen de la priorización}

\begin{center}\small
\begin{tabular}{|l|c|c|c|c|c|}
\hline
\textbf{Categoría} & \textbf{Must} & \textbf{Should} & \textbf{Could} & \textbf{Won\textquotesingle t} & \textbf{Total} \\ \hline
(a) Obligatorio del cliente & 9 & 4 & 3 & --- & 16 \\ \hline
(b) Núcleo de jugabilidad & 8 & 1 & 1 & --- & 10 \\ \hline
(c) Agregado del equipo & --- & --- & --- & 23 & 23 \\ \hline
\textbf{Total} & \textbf{17} & \textbf{5} & \textbf{4} & \textbf{23} & \textbf{49} \\ \hline
\end{tabular}
\end{center}

Los 26 casos \textit{Must}, \textit{Should} y \textit{Could} forman el alcance de las catorce semanas; los 23 de fase posterior quedan documentados para después. Si el plazo aprieta, se recorta primero por la fila \textit{Could} y después por la \textit{Should}; nunca por un \textit{Must}.
